% Challenge 2: the conjugation caveat and the nine-qubit solution.
\begin{frame}{RSP Challenge 2: Complex Conjugation}
  \centering
  \setbeamercolor{itemize item}{fg=CVQCTeal}
  \begin{tikzpicture}[
      >={Latex[length=2.5mm]},
      protocol/.style={draw=CVQCInk,fill=CVQCPale,thick,
        rounded corners=2mm,minimum width=1.7cm,minimum height=0.80cm,
        inner sep=3pt,font=\large,text=CVQCInk},
      output/.style={inner sep=0pt,font=\large,text=CVQCInk},
      bank/.style={output,font=\footnotesize},
      flow/.style={->,very thick,draw=CVQCTeal}
    ]
    \path[use as bounding box] (-6.6,-3.15) rectangle (6.6,3.05);

    \foreach \x in {-5.1,-2.5,0.1,5.1} {
      \node[protocol] at (\x,2.10) {$\pi_{\mathrm{RSP}}$};
      \draw[flow] (\x,1.70) -- (\x,1.22);
    }
    \node[font=\large,text=CVQCMuted] at (2.7,2.10) {$\cdots$};
    \node[font=\large,text=CVQCMuted] at (2.7,0.88) {$\cdots$};

    \only<1>{
      \node[output] at (-5.1,0.88) {$\ket{+_{\theta_1}}$};
      \node[output,text=CVQCCoral] at (-2.5,0.88)
        {$\overline{\ket{+_{\theta_2}}}$};
      \node[output] at (0.1,0.88) {$\ket{+_{\theta_3}}$};
      \node[output,text=CVQCCoral] at (5.1,0.88)
        {$\overline{\ket{+_{\theta_K}}}$};
    }
    \only<2>{
      \node[bank] at (-5.1,0.88)
        {$\bigl(\ket{+_{\theta_{1,j}}}\bigr)_{j=1}^{9}$};
      \node[bank,text=CVQCCoral] at (-2.5,0.88)
        {$\bigl(\overline{\ket{+_{\theta_{2,j}}}}\bigr)_{j=1}^{9}$};
      \node[bank] at (0.1,0.88)
        {$\bigl(\ket{+_{\theta_{3,j}}}\bigr)_{j=1}^{9}$};
      \node[bank,text=CVQCCoral] at (5.1,0.88)
        {$\bigl(\overline{\ket{+_{\theta_{K,j}}}}\bigr)_{j=1}^{9}$};
    }


    \node[anchor=north west,inner sep=0pt,text=CVQCInk] at (-6.1,0.30) {%
      \begin{minipage}{12.2cm}
        \raggedright
        \begin{itemize}
          \setlength{\itemsep}{0.65em}
          \item<1-> The $1/\mathrm{poly}$ guarantee is only up to
            complex conjugation.
          \item<1-> Complex-conjugation attacks are undetectable.
          \item<2-> \textbf{Solution:} Prepare 9 qubits,
            check 8 neighbors for consistency.
        \end{itemize}
      \end{minipage}%
    };
  \end{tikzpicture}

  \note<1>{
    \footnotesize
    The extracted self-test identifies the requested state up to complex
    conjugation, with one flag shared by all nine qubits of a call.
    Separate calls can carry different flags, correlated or coherent with
    the prover's memory. The barred outputs illustrate disagreement, not
    independently sampled classical error bits or an intrinsically invalid
    conjugated state. A single global choice is allowed.
    For phase states, conjugation reverses the phase, with the convention
    used on earlier slides. The self-test alone cannot resolve this
    ambiguity; the checks on the next build control disagreement across
    calls. They do not detect or forbid a common global conjugation.
    The next build explains how the auxiliary qubits test agreement.
    Source: HN extraction and the paper's conjugation-consistency analysis.
  }

  \note<2>{
    \scriptsize
    Each actual call supplies nine qubits: one computation input and
    eight auxiliary qubits sharing its conjugation flag. The displayed
    nine phase states are schematic; the full preparation set also
    includes computational-basis states.
    Assign calls to vertices of a degree-eight expander. Each edge uses
    one auxiliary qubit at each endpoint. The prover reports a Bell-basis
    measurement of this pair, and the verifier checks the outcome against
    its private preparation descriptions. The graph is chosen so that a
    large disagreeing set has many edges crossing to its complement.
    Together with the RSP tests, the analysis bounds contributions with
    too many inconsistent flags, accounting for erroneous preparations.
    In the remaining expansion, fewer than t flags disagree with a
    majority choice. A signed Pauli expansion handles that minority with
    a controlled coefficient cost. It does not physically repair the
    prepared qubits or force every flag to zero.
    For a globally conjugated choice, conjugate the entire prover
    strategy as well. The final classical probabilities are unchanged.
    We then apply the same MBQC, distillation, and evaluation analysis.
    Source: consistency graph, conjugation alignment, and composition
    sections of the submitted paper.
  }
\end{frame}
